\subsection{Area, displacement, and work}

\begin{examplebox}[Area between a graph and the axis]
For $f(x)=2x$ on $[0,3]$,
\[
\int_0^3 2x\,dx=9.
\]
Since the function is non-negative, this is the geometric area under the graph.
\end{examplebox}

\begin{examplebox}[Displacement from velocity]
If
\[
v(t)=2t,
\]
then the displacement from $t=0$ to $t=3$ is
\[
\int_0^3 v(t)\,dt=9.
\]
Because $v(t)\ge0$ on this interval, displacement and distance travelled agree.
\end{examplebox}

\begin{remarkbox}
When velocity changes sign, the integral gives signed displacement. Total distance requires splitting the interval where the direction of motion changes.
\end{remarkbox}

\begin{examplebox}[Work done by a simple variable force]
Suppose a force depends on position according to
\[
F(x)=3x.
\]
The work done from $x=0$ to $x=2$ is
\[
W=\int_0^2 3x\,dx=6.
\]
\end{examplebox}

\begin{interpretationbox}[Units reveal the meaning]
Velocity integrated over time gives distance units. Force integrated over displacement gives work units. The integral combines a rate or density with the variable over which it accumulates.
\end{interpretationbox}
